% 分块：c01（源文件 QuantumWorkTry.tex 第 173--419 行）
% 所属章节：第 1 节 Introduction（导言，含 1.1 / 1.2 / 1.3 三个小节）
% 说明：中文译注稿；行内与行间公式、\label、\ref、\eqref、\cite 及脚注均与原文逐字对应

\section{导言（Introduction）}\label{intro}

\begin{keytakeaways}
\begin{itemize}[leftmargin=1.6em,itemsep=0.2em,topsep=0.2em]
  \item \textbf{物理图像}：SPT 相在无边的紧致空间上只有唯一基态，本身并无通常意义上的拓扑序；但一旦时空带边，体内可逆拓扑场论的相位就无法单独定义，必须由边界上的无能隙费米子以“反常流入”的方式补偿。整数量子霍尔效应的手征边缘模是这一图像的原型。
  \item \textbf{数学工具}：费米子路径积分的指标理论------实表示对应狄拉克算子的模 2 指标，赝实表示对应普通狄拉克指标，复表示对应 APS 指标不变量与 Dai--Freed 定理；自由费米子相的 K 理论分类与配边不变量则提供了统一的框架。
  \item \textbf{核心结论}：按欧几里得号差下表示的类型（赝实、实、复），边界费米子分为三类；3 维拓扑超导体的反常只探测整数值不变量模 16 的约化，与计入相互作用后的 16 类分类相吻合；3 维拓扑绝缘体则给出预期的二值不变量。
\end{itemize}
\end{keytakeaways}

\subsection{本文的目标（Goal Of This Paper）}\label{goal}

对称保护拓扑（symmetry-protected topological，SPT）玻色物相可以用反常（anomaly）来刻画 \cite{WenEtAl}。这一陈述的含义如下。

% [Tutor Note: 原文称配分函数“模为 1”是欧几里得号差下幺正性的推论：可逆（invertible）拓扑场论在每条闭时空上的配分函数取值于单位圆，故可写成指数形式。注意相位只在模 $2\pi$ 意义下有定义；所谓“扣除局部的非普适项”，是指减去与度规、背景规范场相关的局部抵消项（宇宙学常数型项、曲率多项式等），否则短距离自由度会污染相位，使其不能成为拓扑不变量。]
一个 $d$ 维 SPT 相在体内有能隙，且没有通常意义上的拓扑序：在无边的紧致 $d$ 维空间流形 $M$ 上，它具有唯一的基态。但这样的理论并不完全平庸。在维数为 $D=d+1$ 的紧致时空流形 $X$ 上，在大体积极限下，一般地当存在适当的背景规范场时，并在扣除局部的非普适项之后，这类理论的配分函数是一个模为 1 的复数，$Z_X=e^{\i\Phi}$。这里 $e^{\i\Phi}$ 是一个“可逆（invertible）”的拓扑量子场论（topological quantum field theory，TQFT）$\mfT$ 的配分函数，其逆理论具有复共轭的配分函数 $Z_X^{-1}=e^{-\i\Phi}$。

在这种情况下，如果空间流形 $M$（从而时空 $X$）带有边界，则通常无法以一种令人满意的方式------即与一切预期对称性以及幺正性等物理原则都相容的方式------把 $e^{\i\Phi}$ 定义为一个拓扑不变量。要在此情形下定义理论 $\mfT$，就必须为 $M$（以及 $X$）的边界赋予某种附加的自由度。存在若干种可能，但本文所关心的是：沿 $M$ 与 $X$ 的边界存在无能隙的自由度，并且保持理论 $\mfT$ 的对称性。

我们把描述 $X$ 边界上无能隙自由度的理论记为 $\mfB$。该理论具有某种“反常”，并不具备全部所期望的物理性质；因为倘若 $\mfB$ 本身就是一个完全令人满意、保持对称性的理论，那么把它放在 $X$ 的边界 $Y=\partial X$ 上，就根本不能解决最初在 $X$ 上定义理论 $\mfT$ 时所遇到的困难。在 $\partial X$ 上定义理论 $\mfB$ 时的反常，以某种方式抵消了在 $X$ 上定义理论 $\mfT$ 的困难。由此可见，这一情形是------通常是以更抽象的方式------从体内向边界发生反常流入（anomaly inflow）这一观念的范例 \cite{CH}。关于这一情形的通用数学处理，见 \cite{FreedTwo}。这一情形在凝聚态物理中的原型，是带边界的二维流形 $M$ 上的整数量子霍尔效应。在 $M$ 上------或者更确切地说，在三维流形 $X=\R\times M$ 上，其中 $\R$ 参数化时间------存在电磁矢量势 $A$ 的陈--西蒙斯耦合：
\begin{equation}\label{csm}I=k\ICS(A), ~~~\ICS(A)=\frac{e^2}{4\pi}\int_Y\,\d^3x \,\epsilon^{ijk}A_i\partial_jA_k,~~~k\in\Z. \end{equation}
% [Tutor Note: 这是“反常流入”最具体的实现：陈--西蒙斯项在带边流形上的规范变分是一个边界项，其系数恰与边界手征费米子的阿德勒--贝尔--杰克夫反常相反，于是两者相消。另注意 $e^2/4\pi$ 的归一化与 $k\in\Z$ 的量子化：$\ICS(A)$ 只在模 $2\pi\Z$ 意义下有定义，只有 $k$ 为整数时 $e^{\i k\ICS}$ 才良定义，这正是式 (\ref{csm}) 中 $k$ 取整数的原因。]
这一耦合在带边流形上不是规范不变的，相应的“反常”由 $A$ 与在边界上传播的手征边缘模之间的耦合所抵消。（若这些模全都带单位电荷，则其重数为 $k$。）理论 $\mfB$ 描述这些手征边缘模，其反常是手征费米子的阿德勒--贝尔--杰克夫（Adler--Bell--Jackiw）反常的二（时空）维版本 \cite{Adler,BJ}。

近年来研究过的一些最有意思的 SPT 相------其中若干重要情形已在实验中被观测到------是由自由费米子通过能带理论构造出的拓扑绝缘体与拓扑超导体。入门与综述见 \cite{KH,ZQ,MZ}。这些例子中的关键对称性是时间反演对称性 $\sT$，以及与电荷守恒相联系的 $\U(1)$ 对称性。本文的目的是从反常与反常流入的观点来理解这些自由费米子拓扑物相。此前 \cite{RML} 已在这一方向上得到若干结果，其中猜想------正如后文将会清楚的，这一猜想是正确的------考虑整体反常会给出更完整的图像。在分析这一问题时，我们运用指标理论（index theory），把它应用于费米子路径积分。这并不奇怪，因为自由费米子相可以通过 K 理论分类 \cite{Kitaev}（或通过矩阵类分类，二者相互关联 \cite{RSFL}），而指标理论正是在与 K 理论相伴的分析中被发明的 \cite{ASfirst,AS}。

我们应当一开始就说明，从凝聚态物理的角度看，我们的少数考虑是有些异想天开的。我们工作在一个相对论性框架中，其中可以考虑定义在任意 $D$ 维流形 $X$ 上、具有空间边界 $Y$ 的理论。正如大多数反常研究那样，从技术上说取 $X$ 具有欧几里得号差是方便的，我们正是这样做的。之所以能采用相对论性框架，是因为标准拓扑相的边界费米子由带有演生相对论对称性的类狄拉克方程支配；但从凝聚态物理的角度看，依赖这一框架是不自然的。或许可以在对凝聚态物理更自然的哈密顿框架中解释我们的结果，但我们不会尝试这样做。经验似乎表明，有能隙的凝聚态系统常常由完全相对论性的理论支配，但这一点尚未被完全理解。

此外，由于其演生相对论对称性，我们所研究的边界费米子必然具有一种演生对称性；在 $D=4$ 时，这种对称性通常被称为 $\sCPT$（其中 $\sC$ 是电荷共轭，$\sP$ 是宇称或空间反演）。然而，这一表述对于奇数 $D$（偶数空间维数 $d$）并不成立，因为此时 $\sP$ 包含在空间旋转群的连通分支之中，应当代之以一个翻转空间定向的操作。为了采用一种在偶数 $D$ 时与通常的 $\sCPT$ 陈述等价、却又对所有 $D$ 一致适用的语言，我们改而称之为 $\sCRT$ 对称性，其中 $\sR$ 是对某一个空间坐标的反射。$\sCRT$ 对称性在任意维数的相对论理论中总是成立的；因此，尽管它在凝聚态物理中并不很自然，对于一切通常的自由费米子相，它都是 $\partial X$ 上通常的无能隙费米子所具有的演生对称性。这一点将在若干处用到，特别是在第 \ref{welfin} 节中，用来描述与 2 维空间中的拓扑绝缘体相联系的那个拓扑量子场论。
% [Tutor Note: 原文此处疑有笔误：源行 247 写作 “the topological quantum field theory associated a topological insulator”，缺少介词，按语义应为 “associated to/with a topological insulator”。译文按语义补足为“与……相联系”，但保留原文的指代关系。]

% [Tutor Note: 原文强调奇数 $D$ 时不能沿用 $\sCPT$ 而须改用 $\sCRT$，关键在于：当空间维数 $d$ 为偶数时，宇称（空间反演）属于旋转群 $SO(d)$ 的连通分支，不再是与连通分支可交换的独立离散对称性；而“反射一个空间坐标”的操作 $\sR$ 在任意 $d$ 中都有定义。另需注意，洛伦兹号差下 $\sT$ 是反线性的，而这一点在欧几里得号差下看不出来------见本节第一个脚注。]
其实，本导言中谈到拓扑量子场论（TQFT）时，我们一直有些不精确。本文所研究的理论包含费米子，因此需要时空上的自旋结构（spin structure）（在不可定向情形下则是 pin 结构）。于是，所关心的结构类似于 TQFT，只不过要求时空具有一个选定的自旋（或 pin）结构。我们将用 sTQFT（spin topological quantum field theory，自旋拓扑量子场论）这一名称来指称那种类似 TQFT 但包含费米子的理论。关于这类理论在 $D=3$ 陈--西蒙斯理论背景下的早期研究，见 \cite{DW}。sTQFT 中的“s”可以指自旋结构，也可以指 pin 结构，视语境而定。读者或许愿意查阅附录中关于自旋结构与 pin 结构以及相关问题的解释。

\subsection{关于费米子的若干一般性事实（Some Generalities About Fermions）}\label{generalities}

现在我们来解释关于费米子与反常的若干一般性事实，它们将作为本文的有用背景。首先，反常总是只来自无质量费米子；有质量费米子从不贡献反常。一个技术性的解释是：有质量费米子 $\psi$ 所可能具有的任何对称性，也都为泡利--维拉尔斯（Pauli--Villars）正规化子场所具有（后者统计性质相反，满足与 $\psi$ 相同的狄拉克方程，但质量大得多）。因此，有质量费米子总可以以保持其对称性的方式正规化，它们从不贡献反常。

因此，正如直觉所预期的那样，在分析费米子反常时可以忽略有能隙的自由度。事实上，我们还可以忽略可加能隙的自由度------即那些\emph{本可以}具有保持对称性的裸质量（即使实际上并没有）的费米子------因为对这类费米子同样存在保持对称性的正规化子。一个推论是：在任意维数的任何相对论性费米子 $\psi$ 理论中，定义费米子路径积分 $Z_\psi$ 的绝对值 $|Z_\psi|$ 都不成问题，并且与一切对称性和物理原则相容。这是因为，一组共轭费米子 $\t\psi$（它在任何对称性下按 $\psi$ 的复共轭方式变换，其作用量为 $\psi$ 作用量的复共轭）会具有配分函数 $\bar Z_\psi$。费米子的组合系统 $\psi\oplus \t\psi$ 总是可加能隙的。这个组合系统的配分函数为 $Z_\psi \bar Z_\psi=|Z_\psi|^2$。因此，$|Z_\psi|^2$ 没有反常，等价地，$|Z_\psi|$ 也没有反常。

上一段中隐含了这样一点：在欧几里得号差下（与洛伦兹号差相反），费米子未必按相关对称群的实表示或自共轭表示变换（因此 $\psi$ 与 $\t\psi$ 可以按不同方式变换）。例如，整数量子霍尔效应中的边缘模由时空维数为 2 的一个费米子场 $\psi$ 描述，该场具有确定的手征性；在欧几里得号差下，这意味着 $\psi$ 在旋转群 $\SO(2)$ 下按自旋 1/2（或自旋 $-1/2$）变换。边缘模的手征不对称性意味着不存在相应反手征的场，即按 $\psi$ 的复共轭方式变换的场。只有在洛伦兹号差下，费米子才------像玻色场一样------是实的。举例来说，仍以手征边缘态为例：$SO(2)$ 的自旋 $1/2$ 表示中，转角为 $\varphi$ 的旋转作用为 $e^{i\varphi/2}$；在延拓到洛伦兹号差之后，它成为 $SO(1,1)$ 的一个一维表示，其中洛伦兹变换 $\begin{pmatrix}\cosh w & \sinh w\cr \sinh w & \cosh w\end{pmatrix} $ 的作用为 $e^{w/2}$，而这是实的。

一般而言，相对论性费米子按其欧几里得号差下在适当对称群 $K$ 中的变换性质，分为三大类。在 $K$ 中，我们包括旋转群 $SO(D)$、可能的离散对称性（例如\footnote{在欧几里得号差下，一切对称性都以线性方式作用于费米子。$\sT$ 的反线性性质是在解析延拓回洛伦兹号差时才出现的。$\sCRT$ 定理意味着，一个欧几里得对称性可以延拓为洛伦兹号差下两个不同的对称性；例如，欧几里得号差下的对称性 $\sCR$ 延拓到洛伦兹号差后，既可能成为反线性对称性 $\sT$，也可能成为线性对称性 $\sCR$，取决于反射 $\sR$ 是否作用在被解析延拓的那个时空坐标上。} $\sT$ 或 $\sR$），以及规范对称性和/或整体对称性。费米子可以按 (1) 赝实表示（pseudoreal representation）、(2) 实表示（real representation）或 (3) 复表示（complex representation）变换。当然，这些情形的混合也是可能的；例如，一个理论中可能有一些费米子按某一类型的表示变换，另一些按另一类型的表示变换。（此外，一个可约表示可以同时属于不止一种类型；见第 \ref{orch} 节。）我们将在本文中遇到全部三种类型。

% [Tutor Note: 赝实（pseudoreal，亦称四元数型）表示的关键判据是存在不变的反称双线性型 $\omega$：因为 $\omega_{\alpha\beta}$ 反称，$\omega_{\alpha\beta}\psi^\alpha\psi^\beta$ 在交换两个费米子时不改变符号，故可作为裸质量项。这与实表示形成对照：实表示具有不变对称双线性型 $h$，于是 $h_{\alpha\beta}\psi^\alpha\psi^\beta\equiv 0$，裸质量被费米统计禁戒。切勿把 pseudoreal 与 real 混用------二者的区别正是 $\omega$ 反称而 $h$ 对称。]
群 $K$ 的一个不可约表示 $R$ 若具有一个不变的反称双线性型 $\omega$，则称为赝实的。此时，按表示 $R$ 变换的费米子（记作 $\psi^\alpha$，$\alpha=1,\dots,\dim\,R$）可以具有一个 $K$ 不变的裸质量 $\omega_{\alpha\beta}\psi^\alpha\psi^\beta$。（这与费米统计相容，因为 $\omega$ 是反称的。）因此，按赝实表示变换的费米子永远不会贡献反常；严格说来，情形 (1) 与我们的考虑无关。然而，这一类情形有一个变体，我们或可称之为情形 (1$'$)。它指的是这样一个保持 $\sT$ 的理论：其费米子按表示 $R$ 变换，而该表示若在 $K$ 的定义中略去 $\sT$ 便是赝实的，若计入 $\sT$ 则不是。由于在忽略 $\sT$ 时费米子是赝实的，这样的理论总可以以自洽的方式量子化，但 $\sT$ 对称性可能存在反常。这一反常来自涉及费米子路径积分 $Z_\psi$ 符号的问题，如第 \ref{pseudo} 节所解释。基本的例子是 3 维拓扑绝缘体或拓扑超导体的边界费米子，且仅表述在可定向流形上。我们将从这个例子开始，因为其中涉及的拓扑不变量相对而言较为熟悉。此外，我们在 3 维拓扑绝缘体情形下的结果可能特别有意思，因为我们将能够更精确地描述 \cite{ZQH} 中“3 维拓扑绝缘体由等于 $\pi$ 的电磁 $\theta$ 角所刻画”这一陈述的含义。

% [Tutor Note: 情形 (1$'$) 的微妙之处在于：$R$ 在忽略 $\sT$ 的群 $K$ 下是赝实的，因而理论可以无矛盾地量子化、$|Z_\psi|$ 有定义；但把 $\sT$ 计入 $K$ 后 $R$ 不再是赝实的，$\sT$ 不变性可能只差一个符号，即 $Z_\psi$ 的符号反常。这种符号反常不是微扰论能看到的，必须用谱流与指标理论（后文第 \ref{pseudo} 节）来分析，它也正是 $T$ 反常的来源。]
关于情形 (2)：群 $K$ 的不可约表示 $R$ 若具有一个不变的对称双线性型 $h$，则称为实的。对于按这种表示变换的费米子 $\psi$，裸质量被费米统计所禁戒，因为由费米统计可知 $h_{\alpha\beta}\psi^\alpha\psi^\beta =0$。然而，若把费米子谱加倍，即再加入一个同样按 $R$ 的表示变换的场多重态 $\t\psi$，则 $K$ 不变的裸质量 $h_{\alpha\beta}\psi^\alpha\t\psi^\beta$ 就成为可能。谱加倍把费米子路径积分 $Z_\psi$ 替换为其平方 $Z_\psi^2$，因此 $Z_\psi^2$ 无反常；对实费米子而言，反常只能影响 $Z_\psi$ 的符号。我们将考虑的基本例子是 2 维拓扑超导体或拓扑绝缘体的边缘态。

% [Tutor Note: 实表示下 $Z_\psi^2$ 无反常而 $Z_\psi$ 只能相差一个符号，这个符号正是狄拉克算子的模 2 指标（mod 2 index），它是取值为 $0,1$ 的离散量，与连续的 $\eta$ 不变量有本质区别。后文第 \ref{realf} 节会把这一符号写成 $(-1)^\zeta$，其中 $\zeta$ 即模 2 指标。]
最后，关于情形 (3)：$R$ 若不具有任何不变双线性型，则称为复的。此时费米子路径积分 $Z_\psi$ 一般是复的。与其他情形不同，按复表示变换的费米子可以有微扰反常（perturbative anomaly）------即标准的阿德勒--贝尔--杰克夫反常及其在其他维数中的类比。在此语境下，“微扰”反常是指在弱的规范场或引力场中对费米子作量子化时就能看到的那种反常；与之相对的是“整体”或“非微扰”反常，后者在弱背景场中不可见，并涉及整体性的考虑。例如，整数量子霍尔效应中的边缘模按复表示变换，它们确实具有微扰反常（这补偿了陈--西蒙斯耦合 (\ref{csm}) 在带边流形上缺乏规范不变性）。然而，在本文中我们只考虑不存在微扰反常的情形，于是只剩下更微妙的整体反常。

% [Tutor Note: 复表示与实/赝实的本质区别在于：复表示下 $\psi$ 与其复共轭 $\t\psi$ 不等价，因此微扰反常可以出现（即 ABJ 反常及其高维类比）；而本节只讨论微扰反常为零、仅剩整体反常的情形。整体反常用 APS 的 $\eta$ 不变量刻画，但 $\eta$ 本身只在模整数意义下确定，要得到良定义的配分函数必须依靠 Dai--Freed 定理（见下段）。]
% [Tutor Note: 反常的“阶”与取值密切相关：说反常涉及 16 次单位根，等价于说 $Z_\psi$ 取值于 $\mu_{16}=\{e^{2\pi\i n/16}\}$，于是自由费米子的 $\Z$ 不变量只有模 16 的部分能被反常探测；这与计入相互作用后 3 维拓扑超导体分类由 $\Z$ 约化为 $\Z_{16}$ 相符。原文提到 \cite{HCR} 的具体计算只看到 8 次单位根，差别涉及 $\eta$ 不变量的模约定以及整体相位、引力抵消项的处理，详见第 \ref{zeb} 节。]
我们关于复费米子的基本例子，将是一个 3 维拓扑超导体或拓扑绝缘体的边缘模，此时流形可以是不可定向的。特别地，在拓扑超导体情形下，结果发现反常涉及一个 16 次单位根。（此前 \cite{HCR} 曾针对该系统作出具体计算，显示出涉及 8 次单位根的反常；至于反常的阶为何是 16 而不是 8，这一点相当微妙，我们将在第 \ref{zeb} 节中解释。）反常涉及 16 次单位根这一事实意味着：尽管在自由费米子层面拓扑超导体由取值为 $\Z$ 的不变量分类，反常却只能看到该不变量模 16 的约化。事实上，已知 \cite{FCV,WS,MFCV,KitTwo}：计入相互作用后，3 维拓扑超导体的分类由 $\Z$ 约化为 $\Z_{16}$，因此反常恰好探测这 16 个不同的类。
% [Tutor Note: 原文此处疑有笔误：源行 347 写作 “3d topological supeconductors”，漏掉字母 r，应为 “superconductors”。译文按正确拼写处理，特此注明。]
对于 3 维拓扑绝缘体，即便在不可定向流形上，反常也只涉及一个符号，并且反常恰好探测到预期中的 $\Z_2$ 值不变量。

% [Tutor Note: 这一段是作者略去证明的关键处。实表示与赝实表示的反常分别由狄拉克算子的模 2 指标、普通指标拓扑地刻画；而复表示的整体反常则由 APS 的 $\eta$ 不变量表达。Dai--Freed 定理进一步把“边界费米子配分函数的相位”与“体内狄拉克算子的 $\eta$ 不变量”联系起来（粗略地写作 $Z_\psi=|Z_\psi|\,e^{-\i\pi\eta/2}$ 或共轭相位，依号差约定而定）。注意 $\eta$ 只在模 $2$（或模整数）意义下确定，$\eta$ 的这种不确定性恰与引力/规范抵消项的整数移位相消，故指数化后配分函数良定义。]
在欧几里得号差下表示 $R$ 为实的情形，其反常最早在 \cite{Witold,Witoldtwo} 中研究。这类反常可以相当直接地按狄拉克算子的模 2 指标作拓扑理解。赝实费米子的情形可以用类似方式理解，只是用普通的狄拉克指标而非模 2 指标；它最早在 \cite{R,Semenoff,ADM} 中研究。最后，复费米子的整体反常曾在 \cite{Global} 中分析，并用 Atiyah-Patodi-Singer（APS）的 $\eta$ 不变量表示 \cite{APS}。这一关系后来被精炼为一个结果，我们称之为 Dai-Freed 定理 \cite{DF}。Dai-Freed 定理在涉及复费米子的多种问题中都有用，例如杂化弦的世界面路径积分 \cite{Dinstanton}。我们将在第 \ref{comb} 节中用到它。

本工作的一部分动机来自 \cite{KTTW} 的建议：3 维拓扑超导体在 4 维流形上的配分函数应当是某个 $\eta$ 不变量的指数。我们将看到，这一提议可由平凡超导体与拓扑超导体之间相变的标准刻画推出。我们将通过对配分函数的直接计算证明：正如先前工作所提示的，\footnote{正如 \cite{Freed} 中关于式（5.4）的讨论所简略说明的，一般地，取值于 $\U(1)$ 的配边不变量是幺正 TQFT 或 sTQFT 的配分函数。\cite{KTTW} 证明，在本文所考虑的例子里，可能的配边不变量与自由费米子物相一一对应。（在更高维，存在与自由费米子相没有明显联系的配边不变量，因此它们可能与相互作用物相相联系。）在第 \ref{cobo} 节与第 \ref{rpo} 节中，我们就配边不变性如何与幺正性、以及切割与黏合下的行为等物理原则相联系，给出若干很不完整的提示。} 与自由费米子物态相联系的 sTQFT 具有配边不变的配分函数。
\cite{KTTW} 中关于与自由费米子物态相联系的 sTQFT 具有配边不变配分函数的这一提议，可由对体内有能隙费米子的配分函数作直接计算而得到。
% [Tutor Note: 原文此处疑有编辑残留：源行 370--371 先出现一句 “sTQFT's associated to free fermion states of matter have cobordism-invariant partition functions.”，紧接着又以小写 the 起头重复同一命题（“the proposal in [KTTW] that ... follows from direct calculation ...”）。从文气看，前一句应是上句被截断后的重复粘贴，正确的完整句意大概是“正如先前工作所提示的，\cite{KTTW} 关于……的提议可由直接计算得到”。译文保留原文的重复，以便与源码逐句对照；读者可将前一句视为衍文。]
这些体计算与边界费米子反常之间的联系，一般地由 Dai-Freed 定理给出；当无能隙边界费米子是实的或赝实的时，则由更初等的考虑给出。

关于反常的文献浩如烟海，无法在此恰当综述。不过，另有若干论文以与本文相近的精神处理各种问题，它们是：\cite{Bunke}，其中把 \cite{Dinstanton} 关于杂化弦的考虑严格化；\cite{DMW}，其（第 2 节）给出了 11 维流形上 M 理论路径积分的精确定义\footnote{由于实 K 理论的 mod 8 周期性，这个 11 维问题实际上与第 \ref{pseudo} 节所处理的 3 维边界费米子情形有许多共同之处。但细节有所不同，而且反常的抵消不需要来自 12 维的流入。}；\cite{FM}，其中非常精确地处理了该问题的一个更一般版本；\cite{MW}，其（第 5.3 节）在弦论中涉及 D-膜的某个问题里出现了拓扑绝缘体的类比物；以及 \cite{Monnier}，讨论与手征费米子相联系的反常场论。在本文接近完成时，作者注意到在 SPT 相背景下涉及 $\eta$ 不变量的相关工作 \cite{Metlitski}。

% [Tutor Note: 脚注中“mod 8 周期性”指实 K 理论（KO 理论）的 Bott 周期 8：11 维 M 理论路径积分问题与 3 维边界费米子问题落在同一个 mod 8 周期类中，因此可以相互借鉴。但原文特意强调“细节不同、无需 12 维流入”，说明这不是简单的维数平移，不能把 M 理论情形直接等同于 3 维拓扑绝缘体的 $\Z_2$ 情形。]
\subsection{全文结构与一般性说明（Organization and General Remarks）}\label{zorg}

本文内容可概述如下。在第 \ref{pseudo} 节中，我们讨论赝实费米子，或者更确切地说，讨论在忽略 $\sT$ 对称性时为赝实的费米子。在第 \ref{realf} 节中我们讨论实费米子，在第 \ref{complex} 节中讨论复费米子。在第 \ref{maj} 节中，我们简要考虑另一个实费米子的例子：$1+1$ 维中的马约拉纳链 \cite{FK}。总体思路是从简单的例子出发，逐渐引入更复杂的例子。我们在第 \ref{overview} 节中对各种情形作一个概述。

我们尽量使本文正文相对自足，但若干细节放在附录中。关于黎曼几何中旋量的基本定义，见附录 \ref{dirram}。关于低维情形中的细节，见附录 \ref{examples}。4 维（以及 2 维）中的 $\eta$ 不变量在附录 \ref{feta} 中分析。

在本文中，除第 \ref{maj} 节这一唯一例外，我们只考虑满足 $\sT^2=(-1)^F$ 的时间反演对称性（这里 $(-1)^F$ 是按模 2 计算费米子数的算符）。这也意味着，空间反射 $\sR$（在存在 $\U(1)$ 对称性时为 $\sCR$）的平方为 $+1$；并且当我们在不可定向流形上工作时，使用的是 $\Pin^+$ 结构，而不是 $\Pin^-$ 结构。关于这些概念的解释，见附录 \ref{dirram}。在第 \ref{maj} 节中，一切陈述都反转过来：$\sT^2=+1$，$\sR^2=(-1)^F$，并且时空被赋予 $\Pin^-$ 结构。

% [Tutor Note: $\sT^2=(-1)^F$ 与 $\Pin^+$ 的配对是本文的约定，但不同作者对 $\Pin^\pm$ 的命名可能恰好相反（有的把 $\sT^2=(-1)^F$ 对应的结构记作 pin$_-$）。核对文献时务必以 $\sT^2$ 的符号本身为准，而不要只看名称；第 \ref{maj} 节讨论马约拉纳链时全部约定反转（$\sT^2=+1$，$\Pin^-$），正是最容易混淆之处。]
我们经常引入泡利--维拉尔斯正规化子场，以辅助精确地定义费米子的连续统理论。此时正规化子质量的符号是一个重要的变量。在凝聚态物理中，不会有正规化子，而会有一个物理截断。例如，在能带理论中，无能隙的相对论性费米子出现在能带中的某个（或某些）特定点处，而能带其余部分的拓扑所起的作用，类似于泡利--维拉尔斯方法中正规化子质量符号的作用。

% [Tutor Note: 正规化子质量的符号绝非无关紧要的约定：翻转符号相当于把 $\eta$ 不变量平移一个整数，从而改变配分函数的整体相位；在能带图像中，这个符号对应能带其余部分（即“真空”）的拓扑类。原文后文反复用这一点把自由费米子相与配边不变量联系起来，若不抓住“质量符号＝真空拓扑”这一对应，后续关于谱流与 $\theta=\pi$ 的论证会显得突兀。]
我们通常用配分函数来表述论证。一个物理理论远不止其配分函数，但在自由费米子的情形下，只要配分函数有令人满意的定义，关联函数就总有令人满意的定义。因此，把注意力集中在配分函数上，是判断一个理论是否良定义的一种便利的简写方式。偶尔配分函数会为零，这时就必须转而考虑路径积分测度。
